\section{The continuum field and the uniform lattice inputs}
\label{sec:inputs}

We first specify the construction used in this paper and the estimates
we retain from it.  The finite-volume estimates are essential: a gap in
the continuum Hamiltonian alone would not control how many transfer
states can approach its lower edge.

Let $\sigma$ be normalized area measure on $S^2$.  At cutoff $N$, the
lattice spacing is $a=a_N=L^{-N}$, and the probability on a periodic
nearest-neighbor lattice $\Lambda$ is
\begin{equation}
 d\mu_{N,\Lambda}(q)
 =Z_{N,\Lambda}^{-1}
   \exp\!\left(\beta_N\sum_{\{x,y\}\in E(\Lambda)}q_x\cdot q_y\right)
   \prod_{x\in\Lambda}d\sigma(q_x).
 \label{inp:lattice}
\end{equation}
Each bond occurs once; there is no external field.  We use the fixed
prescription of \cite[Theorem~1.2]{OpenAI-O3}: $L$ is a fixed dyadic
integer, the terminal coupling $H$ is sufficiently large, and
\[
 \beta_N=H+\frac{\log L}{2\pi}N+O(\log(N+1)),\qquad B_N>0.
\]
The positive field normalization $B_N$ is the one in that prescription.
For a real test $f\in C_c^\infty(\R^2)$ and a component $i$, put
\begin{equation}
 X_f^i=B_Na^2\sum_x f(x)q_x^i.
 \label{inp:field}
\end{equation}
Coordinates $x$ and all torus lengths are physical.  Torus tests are
periodized from their indicated local charts.  The cutoff plane state
is the limit of periodic square tori, as in
\cite[Proposition~2.1]{OpenAI-O3}; no boundary spins are fixed.
For time reflection, write $(\theta f)(t,b)=f(-t,b)$.

Here are the precise construction results needed below.  They hold for
this same trajectory, normalization, and periodic state.

\begin{proposition}[Inputs from the continuum construction]
\label{inp:contracts}
There are integers $K_0,M_0>0$ and constants $c,C>0$ with the following
properties.  Write $w_k=M_0 2^k$, and let $Z_N(t,w)$ be the partition
function on a torus of time length $t$ and spatial circumference $w$.
\begin{enumerate}[label=\textup{(\roman*)}]
\item In the cutoff plane, all joint moments of the variables
\eqref{inp:field} converge to those of the Euclidean field $\phi$.
The limiting Schwinger distributions are tempered, Euclidean invariant,
and $O(3)$ covariant.  Their Osterwalder--Schrader reconstruction is the
Wightman field of Theorem~\ref{intro:main}, with unique vacuum $\Omega$
and
\[
 H_{\mathrm{phys}}|_{\Omega^\perp}\ge m>0.
\]
For some real test $f$ supported strictly in positive time, the
one-field reflection norm
$\E[\phi^1(\theta f)\phi^1(f)]$ is strictly positive.
\item For all $N\ge K_0$ and $k\ge0$,
\begin{equation}
 0\le \Delta_N(w_k)
 :=4\log Z_N(w_k,w_k)-\log Z_N(2w_k,2w_k)
 \le C e^{-c w_k}.
 \label{inp:defect}
\end{equation}
\item On these square tori and on the cutoff plane, the component
moment measures
\[
 S_{r,N}^{\boldsymbol i}
 =(B_Na^2)^r\sum_{x_1,\ldots,x_r}
     \E\!\left[\prod_{j=1}^r q_{x_j}^{i_j}\right]
     \delta_{(x_1,\ldots,x_r)}
\]
are nonnegative.  For arbitrary translated unit boxes $Q_1,\ldots,Q_r$,
\begin{equation}
 S_{r,N}^{\boldsymbol i}(Q_1\times\cdots\times Q_r)
 \le C^r r!,
 \label{inp:majorant}
\end{equation}
uniformly in $N\ge K_0$, the square torus, the box positions, and the
component indices.
\end{enumerate}
\end{proposition}

For completeness, the source locators distinguish the different
strengths of these inputs.  Part~(i) uses
\cite[Propositions~10.2, 10.3, 11.3, and 11.4]{OpenAI-O3}; the nonzero
one-field norm is proved in the final paragraph of the proof of
Proposition~11.5.  The reconstruction and separated-time
Euclidean continuation are those of
Osterwalder and Schrader~\cite{OsterwalderSchrader1973,OsterwalderSchrader1975}.
Part~(ii) follows from the uniform reference-box test and rectangular
doubling criterion in \cite[Propositions~8.1 and 8.3]{OpenAI-O3}:
the base defect is at most $2^{-10}$ and the defect after $k$ doublings
is at most $64^{-1}2^{-2^k}$.  Their application to the constructed
trajectory is made in \cite[Theorem~8.4]{OpenAI-O3}.  Part~(iii) is
\cite[Proposition~10.1]{OpenAI-O3}.  These statements retain the
fixed parameter choices of that construction, including sufficiently
large $H$.

\subsection{The bottom of the two-point mass support}

Let $\rho$ be the positive spectral measure in
\eqref{intro:spectral}, and define
\[
 \mathcal M=\{\sqrt{s}:s\in\supp\rho\},\qquad m_*=\inf\mathcal M.
\]
The following consequences do not presume that $\rho$ has an atom.

\begin{lemma}[Positive massive Euclidean kernel]
\label{inp:kernel}
One has $0<m_*<\infty$.  At separated Euclidean points, the
like-component two-point function is the continuous kernel
\begin{equation}
 C(t,b)=\int\rho(d\mu^2)\int_{\R}\frac{dp}{2\pi}
          \frac{e^{ipb-tE(p,\mu)}}{2E(p,\mu)},
 \qquad E(p,\mu)=\sqrt{p^2+\mu^2},\quad t>0.
 \label{inp:euclidean}
\end{equation}
It is strictly positive at every nonzero Euclidean separation, and
there is a constant $C_0$ such that
\begin{equation}
 0<C(t,b)\le C_0e^{-m_*\sqrt{t^2+b^2}}
 \quad\text{when }\sqrt{t^2+b^2}\ge1.
 \label{inp:decay}
\end{equation}
\end{lemma}

\begin{proof}
Internal symmetry gives $\E\phi^i(f)=0$, so one-field vectors lie in
$\Omega^\perp$.  Their positive energy-momentum spectral measures
therefore put no weight at energies in $[0,m)$.  If $\supp\rho$ met
$(0,m^2)$, some compact interval $J\subset(0,m^2)$ would have positive
$\rho$-measure.  At sufficiently small spatial momenta, all shells
with mass squared in $J$ have energy below $m$, a contradiction.
Massless spectral weight would likewise give weight at arbitrarily
small positive energies.  Thus $\mathcal M\subset[m,\infty)$.
The nonzero one-field norm in Proposition~\ref{inp:contracts}(i)
implies $\rho\ne0$, so $m_*<\infty$.

The K\"all\'en--Lehmann representation
\cite{Kallen1952,Lehmann1954} and separated-time Euclidean continuation
give \eqref{inp:euclidean}.  To check convergence without an assumption
on the nature of the mass support, recall that a positive tempered
Fourier measure has polynomial growth.  Restricting its spatial
momentum to $|p|\le1$ shows that $\rho([m^2,R^2])$ also has at most
polynomial growth: the mass-shell factor $1/(2E(p,\mu))$ on this
region is bounded below by a constant times $(1+R)^{-1}$.
Exponential damping therefore makes \eqref{inp:euclidean} and its
derivatives absolutely convergent locally for $t>0$.

Euclidean invariance rotates a separation of length $r>0$ to $(r,0)$.
At that point every integrand is positive and $\rho\ne0$.
For $r\ge1$, factor $e^{-m_*r}$ out of the integral and use
\[
 e^{-r(E-m_*)}\le e^{-(E-m_*)}.
\]
The remaining integral is finite by the preceding growth bound.
This proves \eqref{inp:decay}, as well as continuity away from
coincidence.  All uses of the kernel in this paper are off the
coincidence diagonal.
\end{proof}

\subsection{Uniform tails for smeared spins}

The moment input lets us compare field correlations by first comparing
bounded observables.  Its uniformity under separate translations is
important when two tests move far apart.

\begin{lemma}[Moment and tail bounds]
\label{inp:tails}
Fix finitely many real component test shapes of bounded support.
Their arbitrary translates, in the cutoff plane or in the square
tori of Proposition~\ref{inp:contracts}, satisfy
\begin{equation}
 \E|X_f^i|^r\le A^r r!,\qquad
 \E e^{\kappa|X_f^i|}\le C
 \label{inp:tails-bound}
\end{equation}
for some $A,C,\kappa>0$ independent of the cutoff, volume, and
translations.  The constants can also be chosen uniformly for test
families of bounded supremum and support diameter.
For each fixed $r$, tests supported in a common bounded set obey
\begin{equation}
 \|X_f^i-X_g^i\|_{L^r}\le C_r\|f-g\|_\infty,
 \label{inp:test-difference}
\end{equation}
with the same uniformity under translations of that set.
\end{lemma}

\begin{proof}
Cover the support of each test by a bounded number of unit boxes.
For an even moment, nonnegativity of the component moment measure
allows us to replace all test factors by their absolute values and
apply \eqref{inp:majorant}.  This gives
$\E|X_f^i|^{2r}\le A^{2r}(2r)!$.
Cauchy--Schwarz gives the odd absolute moments, after increasing $A$.
Summing the exponential series at any $\kappa<A^{-1}$ proves
\eqref{inp:tails-bound}.  Applying the same argument to $f-g$, with
its supremum factored out, proves \eqref{inp:test-difference}.
The box bound is independent of the box positions, so none of these
constants depends on translations.  Products of any fixed number of
such variables are controlled by H\"older's inequality, even when the
tests are translated independently.
\end{proof}
